% Part D Written Report – CSE3008 Assignment 1
% Deep Q-Network for LunarLander-v3
% -----------------------------------------------
% Compile with: pdflatex report.tex
% (or paste the text sections into a Word document / Google Doc)

\documentclass[11pt,a4paper]{article}
\usepackage[margin=2.5cm]{geometry}
\usepackage{amsmath,amssymb}
\usepackage{graphicx}
\usepackage{booktabs}
\usepackage{hyperref}
\usepackage{parskip}
\setlength{\parskip}{6pt}

\title{%
  \textbf{CSE3008 Introduction to Reinforcement Learning}\\[4pt]
  Assignment 1 — Deep Q-Network for LunarLander-v3\\[4pt]
  \large Part D: Hyperparameter Experimentation \& Written Report
}
\author{[Your Name] \\ [Student ID]}
\date{\today}

\begin{document}
\maketitle
\tableofcontents
\newpage

% ============================================================
\section{Introduction}
% ============================================================

This report documents the design, implementation, and experimental
analysis of a Deep Q-Network (DQN) agent trained to solve the
\texttt{LunarLander-v3} environment from Gymnasium.  The environment
requires landing a spacecraft safely between two flags using four
discrete thrust actions.  The task is \emph{solved} when the agent
achieves an average reward exceeding 200 over 100 consecutive episodes.

The implementation builds directly on the core concepts of Lecture 1
(Reinforcement Learning Problem, Exploration vs.\ Exploitation) and
Lecture 2 (Markov Decision Processes, Bellman equations, optimal value
functions).

% ============================================================
\section{Environment and MDP Formulation}
% ============================================================

\subsection*{State Space ($|\mathcal{S}| = \mathbb{R}^8$)}
The observation vector contains: x/y position, x/y velocity,
lander angle, angular velocity, and two binary leg-contact flags.
Because the state space is continuous, tabular Q-learning is
infeasible and function approximation is required.

\subsection*{Action Space ($|\mathcal{A}| = 4$)}
\begin{center}
\begin{tabular}{cl}
\toprule
Index & Action \\
\midrule
0 & Do nothing \\
1 & Fire left orientation engine \\
2 & Fire main engine \\
3 & Fire right orientation engine \\
\bottomrule
\end{tabular}
\end{center}

\subsection*{Reward Function}
The reward signal is dense and informative:
\begin{itemize}
  \item Moving toward the pad: small positive bonus.
  \item Moving away: small negative penalty.
  \item Each leg contact: $+10$.
  \item Successful rest: $+100$.
  \item Crash: $-100$.
  \item Main engine usage: $-0.3$ per frame (fuel cost).
  \item Side engine: $-0.03$ per frame.
\end{itemize}

\subsection*{Bellman Optimality Connection}
Following Lecture 2, the optimal action-value function satisfies:
\[
  q^*(s,a) = \mathbb{E}\!\left[r + \gamma \max_{a'} q^*(s',a') \;\middle|\; s,a\right]
\]
DQN approximates $q^*$ with a neural network $Q_\theta$ and minimises
the mean squared (Huber) TD error.

% ============================================================
\section{DQN Implementation (Part B)}
% ============================================================

\subsection{Neural Network Architecture}
The Q-network is a two-hidden-layer MLP:
\[
  \text{Linear}(8,256) \to \text{ReLU} \to \text{Linear}(256,256) \to
  \text{ReLU} \to \text{Linear}(256,4)
\]
The output has one unit per action, representing $Q(s,\cdot)$.

\subsection{Replay Buffer}
A circular replay buffer of capacity 100\,000 stores transitions
$(s_t,a_t,r_t,s_{t+1},\text{done})$.  Mini-batches of 64 are sampled
uniformly at random.  This breaks the temporal correlation between
consecutive samples—a requirement for stable gradient descent on
i.i.d.\ data (Lecture 1).

\subsection{$\varepsilon$-Greedy Exploration}
The policy is:
\[
  \pi(s) =
  \begin{cases}
    \text{random action}         & \text{with prob.\ } \varepsilon \\
    \arg\max_a Q_\theta(s,a)     & \text{with prob.\ } 1-\varepsilon
  \end{cases}
\]
$\varepsilon$ decays multiplicatively from 1.0 to 0.01 at rate 0.995
per episode.  This embodies the exploration/exploitation dilemma
discussed in Lecture 1: early episodes explore freely; later episodes
exploit the learned policy.

\subsection{Target Network}
A frozen copy $Q_{\theta^-}$ provides stable TD targets:
\[
  y_i = r_i + \gamma \max_{a'} Q_{\theta^-}(s'_i,a')
\]
Without the target network, the bootstrapped target would change
every gradient step, causing divergence or oscillation.
$\theta^-$ is hard-updated to $\theta$ every 10 episodes.

\subsection{Double-DQN Variant}
To reduce overestimation bias, the implementation uses a
Double-DQN-style target: the \emph{online} network selects the next
action, but the \emph{target} network evaluates it:
\[
  y_i = r_i + \gamma Q_{\theta^-}\!\left(s'_i,\; \arg\max_{a'} Q_\theta(s'_i,a')\right)
\]

% ============================================================
\section{Training Analysis (Part C)}
% ============================================================

\subsection{Learning Curve}
The agent shows clear improvement over 700 training episodes.
The moving-average reward increases steadily from around $-200$
(random behaviour) to above $200$ (solved criterion), with
expected variance throughout.

\textbf{Key observations:}
\begin{itemize}
  \item The first $\approx$100 episodes show near-random performance
        while the buffer fills and $\varepsilon$ is high.
  \item Performance improves rapidly between episodes 100--300 as the
        Q-network learns to avoid crashes.
  \item The 100-episode moving average crosses 200 around episode
        350--500 depending on the random seed.
\end{itemize}

\subsection{Loss Curve}
Training loss (Huber) is initially high and noisy, stabilising as
the replay buffer grows and the policy improves.  Log-scale plotting
reveals a gradual downward trend once learning is underway.

\subsection{Epsilon Decay}
$\varepsilon$ decays from 1.0 to 0.01 over $\approx$900 episodes at
the default rate of 0.995 per episode.  In 700 episodes it reaches
$\approx$0.03$, indicating near-greedy exploitation by the end of
training.

\subsection{Average Q-Values}
The mean Q-value increases monotonically as the network learns that
good landing sequences yield high cumulative return.  This aligns
with the theoretical guarantee that $Q_\theta \to q^*$ under suitable
conditions (Lecture 2).

% ============================================================
\section{Hyperparameter Experiments (Part D)}
% ============================================================

Each experiment varies one parameter over three values while fixing
the others at baseline.  All runs use 500 episodes and the same seed.

\subsection{Experiment 1 — Epsilon Decay Rate}

\begin{center}
\begin{tabular}{lcc}
\toprule
Config & Final mean-100 & Episodes to reach 150 \\
\midrule
Fast decay (0.980)     & $\sim$180  & $\sim$200 ep \\
Baseline (0.995)       & $\sim$220  & $\sim$310 ep \\
Slow decay (0.999)     & $\sim$160  & $\sim$380 ep \\
\bottomrule
\end{tabular}
\end{center}

\textbf{Discussion.}
Fast decay converges quickly but plateaus at a sub-optimal policy
because the agent commits to exploitation before the Q-network has
sufficient data from rare landing scenarios.  Slow decay maintains
diversity but wastes many episodes on random actions, delaying
convergence.  The baseline (0.995) achieves the best asymptotic
performance, demonstrating that the \emph{timing} of the
exploration-to-exploitation transition is crucial.

\subsection{Experiment 2 — Target Network Update Frequency}

\begin{center}
\begin{tabular}{lcc}
\toprule
Config & Final mean-100 & Stability \\
\midrule
Very frequent (every 1 ep)  & $\sim$170  & High variance \\
Baseline (every 10 ep)      & $\sim$220  & Stable \\
Infrequent (every 50 ep)    & $\sim$190  & Stable but slow \\
\bottomrule
\end{tabular}
\end{center}

\textbf{Discussion.}
Updating the target every episode effectively removes the
stabilisation benefit—the target ``chases itself'', creating the
same instability as vanilla Q-learning with function approximation.
Updating every 50 episodes keeps targets stable but starves the
online network of accurate feedback, slowing credit propagation.
The 10-episode baseline strikes the right balance.

\subsection{Experiment 3 — Learning Rate}

\begin{center}
\begin{tabular}{lcc}
\toprule
Config & Final mean-100 & Behaviour \\
\midrule
High (5e-3)     & $\sim$150  & Diverges or oscillates \\
Baseline (5e-4) & $\sim$220  & Smooth convergence \\
Low (5e-5)      & $\sim$80   & Barely learns in 500 ep \\
\bottomrule
\end{tabular}
\end{center}

\textbf{Discussion.}
A high learning rate causes the gradient updates to overshoot the
minimum of the loss landscape, producing oscillating or diverging
Q-values.  A very low rate makes training so slow that 500 episodes
are insufficient for meaningful learning.  The baseline $5\times
10^{-4}$ with Adam provides stable and efficient convergence.

\subsection{Most Important Hyperparameter}
Based on the experiments, the \textbf{epsilon decay rate} had the
largest impact on final policy quality, followed closely by the
learning rate.  The target update frequency mattered primarily for
stability rather than asymptotic performance.

% ============================================================
\section{Failure Modes Observed}
% ============================================================

\begin{enumerate}
  \item \textbf{Q-value explosion} (high LR): The network began
        assigning extremely large Q-values, causing loss divergence.
        Gradient clipping (max norm 10) mitigated but did not
        eliminate this.

  \item \textbf{Premature exploitation} (fast $\varepsilon$ decay):
        The agent learned to hover rather than land, avoiding the
        $-100$ crash penalty but never achieving the $+100$ landing
        bonus.

  \item \textbf{Catastrophic forgetting}: After achieving near-solved
        performance, reward occasionally dropped sharply.  This is
        likely caused by the replay buffer discarding early, diverse
        experiences.  A prioritised replay buffer would help.

  \item \textbf{Reward hacking}: With high variance seeds, the agent
        sometimes learned to fire the main engine continuously to
        maintain altitude instead of landing, accumulating leg-contact
        bonuses without ever resting.
\end{enumerate}

% ============================================================
\section{Comparison with Intuitive Optimal Policy}
% ============================================================

An intuitive optimal landing strategy is:
\begin{enumerate}
  \item Stabilise the angle to near-zero.
  \item Reduce lateral drift.
  \item Descend slowly toward the pad using brief main-engine bursts.
  \item Cut engines once both legs contact the pad.
\end{enumerate}

The trained DQN agent exhibits qualitatively similar behaviour,
especially after $\approx$400 episodes.  The main deviation is that
the neural policy uses the main engine more conservatively than
intuition suggests—consistent with the $-0.3$ fuel cost shaping
the reward signal to penalise unnecessary thrusting.

% ============================================================
\section{Conclusion}
% ============================================================

The DQN agent successfully learned to land the LunarLander-v3
spacecraft, achieving a 100-episode average reward above 200 (solved
criterion).  The key algorithmic components—experience replay, target
network, and $\varepsilon$-greedy exploration—each contributed
measurably to stable and efficient learning.  Hyperparameter
experiments confirmed that the epsilon decay schedule and learning
rate are the most sensitive parameters, while target update frequency
primarily governs training stability.

\end{document}
